\documentclass{article}
\usepackage{graphicx} % Required for inserting images

\usepackage{listings}
\usepackage{xcolor}
\usepackage{subcaption}

\definecolor{codegreen}{rgb}{0,0.6,0}
\definecolor{codegray}{rgb}{0.5,0.5,0.5}
\definecolor{codepurple}{rgb}{0.58,0,0.82}
\definecolor{backcolour}{rgb}{0.95,0.95,0.92}

\lstdefinestyle{mypython}{
    backgroundcolor=\color{backcolour},
    commentstyle=\color{codegreen},
    keywordstyle=\color{blue},
    numberstyle=\tiny\color{codegray},
    stringstyle=\color{codepurple},
    basicstyle=\ttfamily\footnotesize,
    breakatwhitespace=false,
    breaklines=true,
    captionpos=b,
    keepspaces=true,
    numbers=left,
    numbersep=5pt,
    showspaces=false,
    showstringspaces=false,
    showtabs=false,
    tabsize=2,
    language=Python
}

\usepackage[% 
    colorlinks=true,
    pdfborder={0 0 0},
    linkcolor=black
]{hyperref}
\usepackage[a4paper,top=2cm,bottom=2cm,left=3.5cm,right=3.5cm,marginparwidth=1.75cm]{geometry}
\usepackage[export]{adjustbox}
\usepackage{float}
\usepackage{amssymb}
\usepackage{mathtools}
\usepackage{enumitem}
\usepackage{booktabs}

\usepackage{setspace}

\usepackage{titlesec}

% Format: \titlespacing*{<command>}{<left>}{<before-sep>}{<after-sep>}
\titlespacing*{\section}{0pt}{1.2ex plus 1ex minus .2ex}{0.8ex plus .2ex}
\titlespacing*{\subsection}{0pt}{1ex plus 1ex minus .2ex}{0.5ex plus .2ex}
\titlespacing*{\subsubsection}{0pt}{0.8ex plus 1ex minus .2ex}{0.3ex plus .2ex}

% \setstretch{1} % Togli il commento a questo per un valore personalizzato

\title{
    \includegraphics[scale = 0.250]{logo_usi.png}\\
    \vspace{3cm}
    \textbf{\Huge RRL: PPO VS SAM-PPO}
}
 \author{\textbf{ Manuel Tagliaferi \& Carlo Achille Fiammenghi \& Yash Pareek }}
 \date{
 \Large Adavanced Topic in ML A.Y. 2025/26\\
 Professor: Cesare Alippi}

\begin{document}

\maketitle

\newpage

\tableofcontents

\newpage


\section{Introduction}
This report details our project to reproduce and extend the findings of the ICLR 2025 paper \textit{"Flat Reward in Policy Parameter Space Implies Robust Reinforcement Learning"} by Lee and Yoon. In traditional Reinforcement Learning, agents trained in perfect simulators often fail in the real world because they "overfit" to the exact physics of the simulation. To solve this, the authors propose adding Sharpness-Aware Minimization (SAM) to the standard PPO algorithm. Instead of blindly searching for the absolute highest peak of reward, SAM-PPO looks for a "flat" and stable plateau in the policy parameters. 

Our goal was to train both the standard PPO baseline and the SAM-PPO agent on three MuJoCo robots, and thoroughly test which one is truly ready for the real world. We challenged the agents with four distinct physical perturbations: Action Noise (faulty motors), Mass variations (heavier payloads), Friction changes (slippery or sticky floors), and Sensor Noise (inaccurate observations), culminating in a 3D visual analysis of their reward landscapes.

\section{Methodology}

An important preliminary note is that the original codebase provided by the authors was incomplete: it lacked the \texttt{.yaml} configuration files necessary to set the hyperparameters and launch the training scripts. As a result, our first task was to build these configuration files entirely from scratch by carefully reading the paper's appendix and manually implementing the training pipelines.

We trained both the standard PPO baseline and the SAM-PPO agent on three continuous control tasks from the MuJoCo suite:
\begin{itemize}
    \item \texttt{Hopper-v4}
    \item \texttt{Walker2d-v4}
    \item \texttt{HalfCheetah-v4}
\end{itemize}

Each model was trained for $3 \times 10^6$ steps to ensure proper convergence (increasing the default limit of $1 \times 10^6$ steps found in standard tutorials). After training, we evaluated the models against four different types of robustness tests: Action Noise, Mass Variation, Friction Perturbation, and Sensor Noise.

\section{Results and Initial Discrepancies}

\subsection{Hopper and Walker2d: Successful Reproduction}
For the \texttt{Hopper-v4} and \texttt{Walker2d-v4} environments, our reproduction successfully validated the authors' claims. The SAM-PPO agent exhibited significantly less degradation in performance compared to the standard PPO agent as action noise increased. 

\begin{figure}[H]
    \centering
    \includegraphics[width=\textwidth]{"images plotted/comparison_results.png"}
    \caption{Robustness Comparison: Hopper, Walker2d, and HalfCheetah generated by our initial reproduction. Notice how SAM-PPO underperforms in HalfCheetah.}
    \label{fig:comparison1}
\end{figure}

\subsection{The HalfCheetah Discrepancy: A Hyperparameter Issue}
As clearly visible in the third subplot of Figure \ref{fig:comparison}, for the \texttt{HalfCheetah-v4} environment our initial experiments yielded contradictory results: the standard PPO agent outperformed SAM-PPO both in nominal conditions (0 noise) and under noise. Specifically, SAM-PPO converged to an average reward of $\sim 4500$, whereas PPO reached $\sim 6400$.

This plot initially gave us problems, but upon investigation, we discovered it is the result of using an incorrect hyperparameter. We discovered that SAM-PPO is highly sensitive to the neighborhood size parameter, \boldmath{$\rho$}.

\section{Analysis of Hyperparameter Sensitivity}
Due to the absence of an official \texttt{HalfCheetah.yaml} configuration file in the original repository, our initial training run used a \boldmath{$\rho$} that was too high (\boldmath{$\rho = 0.01$}, borrowed from the Hopper configuration). However, a deep dive into Appendix C of the original paper revealed that we must use \boldmath{$\rho = 0.008$} for HalfCheetah.

The use of this excessively large $\rho$ forced the SAM optimizer to search for a wider flat region. In environments without "early termination" (like HalfCheetah, where the robot cannot fall over and fail), this overly conservative constraint prevented the agent from learning a fast running gait.

\section{Corrected HalfCheetah Training and Final Results}
To validate our hypothesis and the authors' findings (specifically Figure 8 of the paper: \textit{Sensitivity of SAM-PPO to $\rho$}), we updated the \texttt{halfcheetah.yaml} configuration to $\rho = 0.008$ and retrained the SAM-PPO agent.

\begin{figure}[H]
    \centering
    \includegraphics[width=\textwidth]{"images plotted/comparison_results2.png"}
    \caption{Corrected Robustness Comparison using $\rho=0.008$ for HalfCheetah. The degradation curve of SAM-PPO is flatter and overtakes standard PPO under severe noise.}
    \label{fig:comparison2}
\end{figure}

The new results perfectly illustrate the classic trade-off between nominal performance and robustness. As seen in the new HalfCheetah subplot (Figure \ref{fig:comparison2}), while standard PPO acts recklessly to achieve a peak nominal reward of $\sim 6400$, it completely collapses to near-zero under severe noise ($\sigma \ge 0.4$). 

In contrast, SAM-PPO learns a more conservative but highly stable gait (nominal reward $\sim 4500$). Its degradation curve is significantly flatter, allowing it to overtake standard PPO at high noise levels ($\sigma \ge 0.35$) and maintain a score of $\sim 1000$ even at maximum perturbation ($\sigma = 0.5$). This empirically validates the paper's claim: flat minima in the parameter space yield policy robustness, preventing catastrophic failures when the environment becomes highly unstable.

\section{Physical Robustness: Mass and Friction Variations}
To further validate the robustness claims of SAM-PPO and provide a more comprehensive analysis, we evaluated the already-trained models against structural and environmental perturbations. Specifically, we tested how the agents handle variations in the robot's body mass and floor friction, without any additional training. 

\begin{figure}[H]
    \centering
    \includegraphics[width=\textwidth]{"images plotted/output3.png"}
    \caption{Mass Robustness Comparison. SAM-PPO exhibits significantly lower variance and avoids the catastrophic failures seen in standard PPO at high mass factors.}
    \label{fig:mass}
\end{figure}

As shown in Figure \ref{fig:mass}, standard PPO acts erratically when the mass of the robot's torso changes. In HalfCheetah, for instance, PPO peaks at the nominal mass (1.0) but fails abruptly when the mass is increased (factors 1.4 and 1.5). Conversely, SAM-PPO retains a smooth, stable, and flat performance dome, remaining highly resilient to weight changes across all three environments.

\begin{figure}[H]
    \centering
    \includegraphics[width=\textwidth]{"images plotted/output4.png"}
    \caption{Friction Robustness Comparison. The flat response of SAM-PPO highlights its invariant performance, compared to the severe oscillations of standard PPO.}
    \label{fig:friction}
\end{figure}

The friction evaluation (Figure \ref{fig:friction}) provides an even more striking result. When altering the floor friction in HalfCheetah, standard PPO exhibits extreme instability, wildly oscillating between $\sim 5000$ and $\sim 6400$ as the environment slightly changes. This high-variance zigzag pattern is the hallmark of a policy that has severely overfitted to the nominal environment. In stark contrast, SAM-PPO maintains a perfectly horizontal line at $\sim 4500$, proving that its learned flat minimum yields an invariant gait, universally stable across both icy and high-friction surfaces.

\section{Sensor Robustness: Observation Noise}
In addition to structural changes and action perturbations, we investigated the resilience of the agents to sensory inaccuracies. By injecting zero-mean Gaussian noise directly into the state observations during evaluation, we simulated an environment with faulty or noisy sensors.

\begin{figure}[H]
    \centering
    \includegraphics[width=\textwidth]{"images plotted/output5.png"}
    \caption{Sensor Noise Robustness. Performance degradation as observation noise increases.}
    \label{fig:sensor_noise}
\end{figure}

As depicted in Figure \ref{fig:sensor_noise}, SAM-PPO demonstrates superior stability even when its "vision" of the environment is distorted. Standard PPO's reliance on a highly specific, sharp minimum makes it extremely sensitive to inaccurate state readings, causing its performance to drop sharply. SAM-PPO's flat policy parameters, however, ensure that slight miscalculations in the state representation do not lead to drastically wrong actions, preserving a much higher average reward.
\section{Joint Perturbations and the Robustness Trade-off}

To visualize the concept of "flat minima" in a truly comprehensive way, we evaluated the \texttt{HalfCheetah} agents on a 2D grid of simultaneous perturbations (varying both mass and friction factors concurrently from 0.6x to 1.4x). 

\begin{figure}[H]
    \centering
    \includegraphics[width=\textwidth]{"images plotted/output6.png"}
    \caption{2D Heatmap of Joint Mass and Friction Perturbations. PPO displays a sharp performance peak surrounded by failure zones, while SAM-PPO maintains a consistent plateau.}
    \label{fig:heatmap}
\end{figure}

The heatmap (Figure \ref{fig:heatmap}) highlights the raw numerical performance across 25 different environmental configurations. To make this topological difference even clearer, we used \texttt{SciPy}'s bivariate spline interpolation to render the resulting high-resolution 3D reward landscapes.

\begin{figure}[H]
    \centering
    \includegraphics[width=\textwidth]{"images plotted/output7.png"}
    \caption{High-Resolution 3D Reward Landscape Comparison. Standard PPO (left) forms a high but extremely fragile peak. SAM-PPO (right) forms a slightly lower but remarkably stable plateau.}
    \label{fig:landscape}
\end{figure}

Figure \ref{fig:landscape} provides the most striking visual confirmation of the paper's core thesis. The Standard PPO landscape (left) features a very high peak at the nominal center (Mass=1.0, Friction=1.0) reaching nearly 6500 reward points. However, this peak is a "sharp minimum": moving slightly in any direction causes the agent to fall off a mathematical cliff, resulting in severe performance drops (deep red valleys).

In stark contrast, SAM-PPO (right) never reaches the 6500 peak. Instead, it deliberately seeks out a "flat minimum" during training, resulting in a broad, stable plateau of $\sim 4500$ points. 

This directly highlights the fundamental \textbf{Robustness-Nominal Performance Trade-off} in Reinforcement Learning: to achieve true robustness, an agent must sacrifice peak nominal performance. A policy optimized purely for a perfect, static environment (like standard PPO) will be fragile and brittle in reality. SAM-PPO trades a marginal loss in theoretical top-speed for the guarantee that the robot will not fail when deployed in an imperfect, unpredictable real world.

\section{Conclusion}
In this project, our primary goal was to verify if the SAM-PPO algorithm is truly more robust than standard PPO. We successfully reproduced the core experiments on three MuJoCo environments: \texttt{Hopper-v4}, \texttt{Walker2d-v4}, and \texttt{HalfCheetah-v4}.

During our work, we faced an interesting challenge: initially, SAM-PPO failed to learn a good gait on HalfCheetah. Instead of assuming the code was wrong, we analyzed the algorithm and discovered its high sensitivity to the neighborhood size $\rho$. By correcting this value from $0.01$ to $0.008$ (as hidden in the paper's appendix), we fixed the issue and obtained excellent results.

Our extensive testing proved the paper's thesis across four different real-world scenarios:
\begin{enumerate}
    \item \textbf{Action Noise:} When adding random noise to the motors, standard PPO crashed to zero, while SAM-PPO degraded slowly and gracefully.
    \item \textbf{Mass Variation:} When making the robot heavier, PPO failed to carry the extra weight, while SAM-PPO kept running almost flawlessly.
    \item \textbf{Friction Variation:} When changing the floor from icy to sticky, PPO showed extreme instability (proof of severe overfitting), whereas SAM-PPO maintained a perfectly flat and invariant performance.
    \item \textbf{Sensor Noise:} When confusing the robot with inaccurate sensor readings, SAM-PPO's flat policy prevented it from making catastrophic decisions, outperforming PPO by a large margin.
\end{enumerate}

Finally, our 3D Joint Robustness Landscape provided the definitive visual proof. This project taught us a fundamental lesson in Reinforcement Learning: optimizing only for the highest possible score in a perfect environment creates fragile, "glass-cannon" policies. SAM-PPO proves that finding "flat" regions in the parameter space---sacrificing a bit of theoretical top-speed---is the absolute key to gaining massive advantages in safety, stability, and real-world deployment.

\end{document}

